%%%%%%%%%%%%%%%%%%%%%%%%%%%%%%%%%%%%%%%%%%%%%%%%%%%%%%%%%%%%%%%%%%%%%%
\section{Objectifs de la première phase}
%%%%%%%%%%%%%%%%%%%%%%%%%%%%%%%%%%%%%%%%%%%%%%%%%%%%%%%%%%%%%%%%%%%%%%

La première phase du projet ALORA a pour objectif principal de concevoir
et de mettre en place l'ensemble des ressources nécessaires au
développement futur de modèles d'intelligence artificielle dédiés au
traitement automatique de la parole en soussou.

Contrairement à une approche centrée directement sur la modélisation, ce
travail privilégie la constitution d'un corpus oral structuré, fiable et
documenté, qui servira de fondation à toutes les expérimentations
ultérieures.

Les objectifs de cette première phase sont les suivants :

\begin{enumerate}

\item Concevoir une plateforme de collecte collaborative permettant
l'enregistrement de données vocales par des locuteurs natifs.

\item Définir une architecture de données inspirée des bonnes pratiques
utilisées dans les corpus de référence, notamment WAXAL, afin de garantir
la cohérence, la traçabilité et la réutilisabilité des données.

\item Construire un premier corpus oral comprenant les enregistrements,
les métadonnées associées ainsi que les informations nécessaires à leur
exploitation.

\item Mettre en place une chaîne de traitement permettant la
transcription, la segmentation et l'annotation progressive des données
collectées.

\item Produire un corpus directement exploitable pour les futures phases
de modélisation, d'analyse linguistique et d'apprentissage automatique.

\end{enumerate}

À l'issue de cette première phase, le principal livrable attendu ne sera
pas uniquement la plateforme de collecte, mais un corpus oral organisé,
documenté et validé, pouvant servir de ressource de référence pour les
travaux de recherche menés dans le cadre du projet ALORA.

\begin{important}

L'objectif de cette phase n'est pas de développer un modèle de
reconnaissance vocale ou de traduction, mais de construire les
fondations scientifiques et techniques sur lesquelles ces modèles
pourront être développés dans les étapes suivantes du projet.

\end{important}